\documentclass[12pt,a4paper]{article}
\usepackage[utf8]{inputenc}
\usepackage[russian]{babel}
\usepackage{amsmath,amssymb,amsfonts}
\usepackage{geometry}
\geometry{top=2cm,bottom=2cm,left=2cm,right=2cm}
\pagestyle{empty}

\begin{document}

\begin{center}
    {\large\bf Листок 1. Принцип Дирихле} \\
    \vspace{2mm}
    {\small Математический кружок}
\end{center}

\vspace{5mm}

\begin{enumerate}
    \item В коробке лежат 10 черных и 10 белых носков. Какое наименьшее количество носков нужно вытащить не глядя, чтобы среди них гарантированно оказалась хотя бы одна пара носков одного цвета?
    
    \item В классе учится 30 человек. Докажите, что найдутся хотя бы два ученика, у которых фамилии начинаются с одной и той же буквы.
    
    \item Докажите, что среди любых $n+1$ целых чисел можно выбрать два, разность которых делится на $n$.
\end{enumerate}

\end{document}
